\section{Implement Lower Bound}
\label{sec:bs-lower-bound}

\problemstatement{We are given a sorted array $\textit{arr}$ of $n$
elements and a value $X$. We must find the \textbf{lower bound} of $X$:
the index of the first element in $\textit{arr}$ that is $\textbf{greater
than or equal to}$ $X$. If no such element exists (every element is smaller
than $X$), the lower bound is $n$ (one past the end of the array).}

\begin{patternbox}
\emph{``Find the first/smallest index such that some condition holds''} --
here, the condition is $\textit{arr}[i] \ge X$ -- is the
\textbf{find-the-boundary} variant of binary search, distinct from the
exact-match search of Section~\ref{sec:bs-find-x}. The keyword is not
``find $X$'' but \textbf{``find the first position where a property
becomes true.''} This phrasing recurs constantly throughout this chapter
(first/last occurrence, search insert position, floor/ceil, and many
``binary search on the answer'' problems), so internalizing the
boundary-search template here pays off many times over.
\end{patternbox}

\subsection*{Understanding the problem carefully}

Define the predicate $P(i) = (\textit{arr}[i] \ge X)$. Because the array is
sorted in increasing order, $P$ is \texttt{false} for some prefix of
indices and \texttt{true} for the remaining suffix, with exactly one
crossover point -- the lower bound is, by definition, that crossover index.
Unlike Section~\ref{sec:bs-find-x}, we are not looking for a single exact
match; we are looking for the \emph{boundary} between ``too small'' and
``big enough,'' which exists regardless of whether $X$ itself is actually
present in the array.

\begin{ideabox}[Candidate Space and Predicate]
The candidate space is indices $0, \dots, n$ (note: $n$ itself is a valid
answer, representing ``no element is large enough''). The predicate is
$P(i) = (\textit{arr}[i] \ge X)$ for $i < n$, and we want the smallest index
where $P$ holds.
\end{ideabox}

\subsection*{Why we keep a separate ``best answer so far'' variable}

In exact-match search, the loop stops the instant we find a match. Here, we
cannot stop early, because finding \emph{one} index where $\textit{arr}[i]
\ge X$ does not tell us it is the \emph{first} such index -- there might be
an even smaller valid index hiding in the left half. So instead of
returning immediately, whenever $\textit{arr}[\textit{mid}] \ge X$ we
record $\textit{mid}$ as a candidate answer and \emph{keep searching to the
left}, hoping to find an even smaller valid index; whenever
$\textit{arr}[\textit{mid}] < X$, we know the entire left portion up to and
including $\textit{mid}$ is too small, so we discard it and search to the
right. This is the general template for every ``find the first/last index
satisfying a monotonic property'' problem: maintain a running best answer,
and keep shrinking the window based on the predicate, rather than stopping
on the first observed success.

\subsection*{Algorithm}

\begin{enumerate}
  \item Initialize $\textit{lo} \gets 0$, $\textit{hi} \gets n-1$,
        $\textit{ans} \gets n$.
  \item While $\textit{lo} \le \textit{hi}$:
  \begin{enumerate}
    \item $\textit{mid} \gets \textit{lo} + (\textit{hi}-\textit{lo})/2$.
    \item If $\textit{arr}[\textit{mid}] \ge X$: record
          $\textit{ans} \gets \textit{mid}$, and search left:
          $\textit{hi} \gets \textit{mid}-1$.
    \item Otherwise: search right: $\textit{lo} \gets \textit{mid}+1$.
  \end{enumerate}
  \item Return \textit{ans}.
\end{enumerate}

\begin{algorithm}[H]
\caption{Lower Bound}
\begin{algorithmic}[1]
\Procedure{LowerBound}{$\textit{arr}, X$}
  \State $\textit{lo} \gets 0$, $\textit{hi} \gets n-1$, $\textit{ans} \gets n$
  \While{$\textit{lo} \le \textit{hi}$}
    \State $\textit{mid} \gets \textit{lo} + (\textit{hi}-\textit{lo})/2$
    \If{$\textit{arr}[\textit{mid}] \ge X$}
      \State $\textit{ans} \gets \textit{mid}$
      \State $\textit{hi} \gets \textit{mid} - 1$
    \Else
      \State $\textit{lo} \gets \textit{mid} + 1$
    \EndIf
  \EndWhile
  \State \Return \textit{ans}
\EndProcedure
\end{algorithmic}
\end{algorithm}

\subsection*{Dry run}

\dryrun{Take $\textit{arr} = [1, 2, 2, 3]$, $X = 2$.}

\begin{itemize}
  \item $\textit{lo}=0,\textit{hi}=3,\textit{ans}=4$: $\textit{mid}=1$,
        $\textit{arr}[1]=2 \ge 2$. $\textit{ans}=1$, $\textit{hi}=0$.
  \item $\textit{lo}=0,\textit{hi}=0$: $\textit{mid}=0$,
        $\textit{arr}[0]=1 \ge 2$? No. $\textit{lo}=1$.
  \item $\textit{lo}=1 > \textit{hi}=0$: loop ends.
\end{itemize}

The lower bound is index $1$ -- the first occurrence of $2$.

Now take $X = 5$ on the same array (larger than every element).

\begin{itemize}
  \item $\textit{lo}=0,\textit{hi}=3,\textit{ans}=4$: $\textit{mid}=1$,
        $\textit{arr}[1]=2 \ge 5$? No. $\textit{lo}=2$.
  \item $\textit{lo}=2,\textit{hi}=3$: $\textit{mid}=2$,
        $\textit{arr}[2]=2 \ge 5$? No. $\textit{lo}=3$.
  \item $\textit{lo}=3,\textit{hi}=3$: $\textit{mid}=3$,
        $\textit{arr}[3]=3 \ge 5$? No. $\textit{lo}=4$.
  \item $\textit{lo}=4 > \textit{hi}=3$: loop ends. $\textit{ans}$ remains
        $4$ (its initial value, never overwritten).
\end{itemize}

The lower bound is $4 = n$, correctly signaling that no element is large
enough.

\subsection*{C++ implementation}

\begin{lstlisting}
#include <bits/stdc++.h>
using namespace std;

int lowerBound(vector<int>& arr, int X) {
    int n = arr.size();
    int lo = 0, hi = n - 1, ans = n;

    while (lo <= hi) {
        int mid = lo + (hi - lo) / 2;
        if (arr[mid] >= X) {
            ans = mid;
            hi = mid - 1;
        } else {
            lo = mid + 1;
        }
    }
    return ans;
}

int main() {
    vector<int> arr = {1, 2, 2, 3};
    cout << "Lower bound of 2: " << lowerBound(arr, 2) << endl;
    cout << "Lower bound of 5: " << lowerBound(arr, 5) << endl;
    return 0;
}
\end{lstlisting}

\begin{complexitybox}
\textbf{Time Complexity.} The search window halves each iteration, giving
\[
  O(\log n).
\]
\textbf{Space Complexity.} Only a constant number of scalar variables are
used, giving $O(1)$ space.
\end{complexitybox}

\begin{takeawaybox}
Whenever a problem asks for the \emph{first} (or \emph{last}) index
satisfying a monotonic condition, do not stop the binary search the moment
you observe the condition holding -- instead, record it as a tentative
answer and keep narrowing the search toward the boundary, in the direction
that could produce an even better (earlier, or later) answer. This
``record and keep narrowing'' template is reused, almost verbatim, in
upper bound, floor/ceil, and first/last occurrence.
\end{takeawaybox}
